\chapter{Geometry and the mathematical definition of manifolds}\label{ch:03}
\begin{insight}{Learning objectives}
Read a manifold definition rather than substitute a visual shape for it; construct charts for familiar examples; distinguish a parameterization, immersion, and embedding; calculate tangent spaces and induced metrics; and identify the statistical assumptions hidden in manifold learning.
\end{insight}

\section{Why one coordinate may require several charts}
A small part of a circular track can be described by one number, distance along the track. A complete circular track cannot be represented faithfully by a single real interval with its usual endpoint structure unless the endpoints are identified. Locally it is one-dimensional, but globally it closes onto itself. A manifold formalizes this separation of local coordinates from global organization.

Before defining manifolds, we need to say what ``nearby'' and ``continuous'' mean without assuming the entire object already has Cartesian coordinates. This is the role of metrics and topology. Geometry then adds quantitative notions such as length and angle.

\section{Metric spaces and topological spaces}
\begin{definition}[Metric]
A metric on a set $S$ is a function $d:S\times S\to[0,\infty)$ such that $d(x,y)=0$ exactly when $x=y$, $d(x,y)=d(y,x)$, and $d(x,z)\leq d(x,y)+d(y,z)$. An open ball is $B_r(x)=\{y:d(x,y)<r\}$, where $r>0$. A set $U$ is open if every $x\in U$ has some ball $B_r(x)$ contained in $U$.
\end{definition}
Open does not mean unbounded, and closed does not mean compact. The open interval $(0,1)$ is bounded, the closed line $\R$ is unbounded, and a set may be both open and closed in its ambient space. A neighborhood of $x$ means a set containing an open set that contains $x$.

\begin{definition}[Topology and continuity]
A topology $\mathcal T$ on $S$ is a collection of subsets called open sets, containing $\varnothing$ and $S$, closed under arbitrary unions and finite intersections. A map $f:S\to T$ between topological spaces is continuous when the inverse image $f^{-1}(U)$ of every open set $U\subseteq T$ is open. A homeomorphism is a bijection that is continuous and whose inverse is continuous.
\end{definition}
Every metric supplies a topology through its open balls. Different metrics can give the same topology while assigning different distances. Topology preserves the structure of continuous neighborhoods, not numerical lengths or angles. A homeomorphism is the appropriate notion of sameness in topology; an isometry, which preserves distances, is a stronger geometric notion.

For a subset $M\subset\R^p$, the \emph{subspace topology} declares $U\subset M$ open exactly when $U=M\cap O$ for some open $O\subset\R^p$. Thus a short arc is open in a circle even though it is not open in the whole plane. The ambient space must always be specified when discussing openness.

Two modest regularity conditions rule out pathological coordinate spaces. \emph{Hausdorff} means distinct points have disjoint open neighborhoods. \emph{Second countable} means a countable collection of open sets forms a basis: every open set is a union of members of that collection. Euclidean spaces and their subspaces have these properties. They ensure that a manifold behaves like a manageable geometric space, rather than an arbitrary collection of local coordinate pictures.

\section{Topological and smooth manifolds}
\begin{definition}[Topological manifold]\label{def:manifold}
A $d$-dimensional topological manifold without boundary is a Hausdorff, second-countable topological space $M$ such that every point has an open neighborhood $U\subset M$ homeomorphic to an open subset $V\subset\R^d$. Such a homeomorphism $\varphi:U\to V$ is a coordinate chart. A collection of charts covering $M$ is an atlas.
\end{definition}
Every part of this definition has a purpose. ``Every point'' rules out a singular junction hidden in an otherwise smooth-looking set. ``Open subset'' rules out endpoints for a manifold without boundary. ``Homeomorphic'' requires both a continuous coordinate map and a continuous reconstruction. ``Locally'' allows the global object to wrap, close, or have holes. Manifolds with boundary instead permit local models in a half-space $\{u_d\geq0\}$; a closed interval is a one-dimensional manifold with boundary, but its endpoints prevent it from being one without boundary.

\begin{example}[What qualifies, and what does not]
A line and a plane are manifolds of dimensions one and two. Every affine $d$-plane is a $d$-manifold. The union of two lines crossing at a point is not a one-dimensional manifold there: removing that point from a sufficiently small neighborhood leaves four branches, whereas removing the center of an interval leaves two. A finite set of distinct observations with the Euclidean subspace topology is a discrete zero-dimensional manifold. Therefore the literal topology of a finite point cloud is not the intended positive-dimensional manifold in a sampling model.
\end{example}

Charts alone allow continuous deformations with corners or irregular derivatives. To differentiate consistently, suppose two charts $\varphi_\alpha:U_\alpha\to V_\alpha$ and $\varphi_\beta:U_\beta\to V_\beta$ overlap. Their transition map
\begin{equation}
\varphi_\beta\circ\varphi_\alpha^{-1}:
\varphi_\alpha(U_\alpha\cap U_\beta)\longrightarrow
\varphi_\beta(U_\alpha\cap U_\beta)
\end{equation}
changes coordinates for the same point. A smooth atlas requires every such transition to have derivatives of all orders. A \emph{smooth manifold} is a topological manifold equipped with a maximal compatible smooth atlas. ``Maximal'' includes every chart compatible with the chosen smooth structure; it prevents an arbitrary choice of two versus three charts from defining different smooth manifolds. A $C^r$ manifold requires only $r$ derivatives. The usual differential-geometric results used here concern smooth or explicitly stated $C^r$ structures \citep{lee2012}.

\begin{misconception}
A smooth-looking two-dimensional visualization does not verify the chart definition. A fitted map can be smooth by construction even when the underlying data include branches, mixed dimensions, outliers, or disconnected components. The mathematical regularity of an estimator and the regularity of the process generating its input are different assertions.
\end{misconception}

\section{Canonical examples and their local coordinates}
\subsection{The circle}
The unit circle is $S^1=\{(x,y)\in\R^2:x^2+y^2=1\}$. The map $\theta\mapsto(\cos\theta,\sin\theta)$ parameterizes it periodically. Restricted to any sufficiently short open interval it is a smooth local coordinate inverse. The notation $\theta\bmod 2\pi$ identifies values differing by integer multiples of $2\pi$. Such an angle is circle-valued, not a globally continuous real-valued coordinate.

A continuous bijection from $[0,2\pi)$ to $S^1$ exists through this formula, but its inverse is discontinuous at the identified direction. It therefore fails to be a homeomorphism. This supplies a concrete reason why injectivity and continuity alone do not guarantee an embedding for noncompact domains.

\subsection{The sphere}
The unit sphere is $S^2=\{(x,y,z):x^2+y^2+z^2=1\}$. Away from the north pole $(0,0,1)$, stereographic coordinates are
\begin{equation}
(u,v)=\left(\frac{x}{1-z},\frac{y}{1-z}\right),\qquad
(x,y,z)=\frac{(2u,2v,u^2+v^2-1)}{1+u^2+v^2}.
\end{equation}
The denominator is positive, so the inverse is smooth on the whole coordinate plane. A second chart from the south pole covers the missing point. Transition maps are smooth wherever both charts apply. Latitude and longitude are useful operational coordinates, but their polar singularities belong to the coordinate system, not to the sphere itself.

\subsection{The torus and a high-dimensional curve}
An abstract torus is $T^2=S^1\times S^1$, described by two independent periodic angles. A familiar embedding in three-dimensional space is
\begin{equation}\label{eq:torus-embedding}
q(\theta,\psi)=\big((R+r\cos\psi)\cos\theta,
(R+r\cos\psi)\sin\theta,r\sin\psi\big),\quad R>r>0.
\end{equation}
The inequality prevents self-intersection and degeneracy in this parameterization. A square with opposite edges identified represents the same abstract topology; drawing it flat does not preserve all distances of the doughnut surface. A torus therefore illustrates the distinction between an abstract manifold and a particular embedding with a particular metric.

For a high-dimensional example, $q(s)=(s,s^2,\sin s,e^s)$ on an open interval is a smooth embedded curve in $\R^4$. Its derivative has first coordinate one, so the derivative never vanishes; its first coordinate also makes the map one-to-one with a continuous inverse. Four observed coordinates and one local degree of freedom are compatible.

\section{Differentiable maps, tangent vectors, and Jacobians}
A map between smooth manifolds is smooth when its expression in overlapping source and target charts is smooth in Euclidean coordinates. This definition makes differentiability independent of the chosen coordinates.

For an embedded manifold $M\subset\R^p$ and a point $x\in M$, a tangent vector is the velocity $\dot c(0)$ of a smooth curve $c$ in $M$ with $c(0)=x$. These velocities form a $d$-dimensional vector space $T_xM$, the tangent space. If $q:V\subset\R^d\to M$ is a regular local parameterization and $q(u)=x$, then
\begin{equation}\label{eq:tangent-param}
T_xM=\{Dq(u)a:a\in\R^d\}.
\end{equation}
The columns of $Dq$ are coordinate tangent vectors. They need not be orthonormal. For an abstract manifold, tangent vectors can be defined as equivalence classes of curves having the same velocity in a chart; the chain rule shows this equivalence does not depend on the chart. This construction avoids assuming an ambient Euclidean space.

A smooth map $f:M\to N$ has a differential $Df_x:T_xM\to T_{f(x)}N$, a linear map sending velocities to velocities. For a curve $c$, $\frac{d}{dt}f(c(t))|_{t=0}=Df_x\dot c(0)$. A Jacobian is the matrix of this coordinate-independent linear map in chosen coordinates.

\begin{example}[Tangents to a sphere]
Differentiate $c(t)^Tc(t)=1$ at $t=0$ to obtain $2x^T\dot c(0)=0$. Thus $T_xS^2=\{v:x^Tv=0\}$, the plane perpendicular to the radius, translated to the origin as a vector space. At $(0,0,1)$ it is the horizontal plane. A local PCA estimate of this plane is an estimate of tangent orientation; it does not reveal whether the full object is a sphere, torus, or an open patch.
\end{example}

\section{Inverse and implicit function theorems}
The rank of a Jacobian measures how many independent infinitesimal changes survive a map. It supplies local conditions for coordinates and constraints.

\begin{theorem}[Inverse function theorem]\label{thm:inverse-function}
Let $f:U\subset\R^d\to\R^d$ be $C^1$, with $U$ open, and suppose $\det Df(x_0)\ne0$. There are neighborhoods of $x_0$ and $f(x_0)$ on which $f$ is a bijection with a $C^1$ inverse. Its inverse derivative is $D(f^{-1})_{f(x)}=[Df_x]^{-1}$. If $f$ is smoother, the local inverse has the corresponding smoothness.
\end{theorem}
The theorem is local: $f(s)=\sin s$ has nonzero derivative near zero but is not globally one-to-one. Nonzero determinant is sufficient for a differentiable local inverse; a continuous inverse can exist when it fails, as for $s\mapsto s^3$. The derivative formula follows by differentiating $f^{-1}\circ f=\operatorname{id}$.

\begin{theorem}[Implicit function and regular level sets]\label{thm:implicit-function}
Let $F:U\subset\R^p\to\R^q$, with $U$ open, be $C^1$, $q\leq p$. If $F(x_0)=0$ and $DF(x_0)$ has rank $q$, then near $x_0$ the zero set can be written, after rearranging coordinates, as the graph of a $C^1$ function of $p-q$ variables. If this rank condition holds at every point of $F^{-1}(0)$, the level set is a $C^1$ embedded manifold of dimension $p-q$, with $T_xM=\ker DF(x)$.
\end{theorem}
\begin{proof}[Proof idea]
Select $q$ coordinates for which the corresponding $q\times q$ derivative block is invertible. Keep the remaining $p-q$ coordinates and replace the selected ones with the $q$ values of $F$. This combined coordinate map has invertible derivative, so the inverse function theorem applies. In the new coordinates the zero set is simply the plane on which the last $q$ coordinates vanish. Differentiating $F(c(t))=0$ yields the tangent-space formula.
\end{proof}
For $F(x,y)=x^2+y^2-1$, the derivative $(2x,2y)$ has rank one on the circle, proving it is a one-dimensional embedded manifold. For $F(x,y)=xy$, the derivative vanishes at the origin; the zero set is crossed axes and is singular there. Rank failure means the theorem is inconclusive, not automatically that the set is singular: $F(x,y)=y^3$ also has vanishing derivative on its zero set, yet that zero set is an ordinary line.

\section{Immersions and embeddings}
\begin{definition}[Immersion and smooth embedding]
A smooth map $q:M^d\to N^p$ is an immersion if $Dq_x$ is injective at every point, equivalently has rank $d$. It is a smooth embedding if it is also a homeomorphism from $M$ onto its image with the subspace topology. An embedded submanifold is the image equipped with this compatible smooth structure.
\end{definition}
An immersion preserves local tangent dimension but may identify distant points. The curve $q(t)=(\sin t,\sin2t)$ on $S^1$ is an immersion because its derivative $(\cos t,2\cos2t)$ never vanishes, yet it crosses itself and is not injective. Even an injective immersion need not be an embedding: an irrational straight-line winding on a torus has a dense image and an inverse that is not continuous in the image topology. Compactness helps: a continuous injective map from a compact space into a Hausdorff space is a homeomorphism onto its image. Thus a smooth injective immersion of a compact manifold is an embedding.

\begin{insight}{The observation map matters}
If a latent state lies on $M$ and observations are $y=g(x)$, the observed image has the same smooth manifold structure only under suitable conditions on $g$, such as a smooth embedding. A projection can merge distinct latent states, collapse dimensions, or create apparent intersections. Recovering the geometry of observed activity is therefore not automatically recovering the geometry of the full neural state.
\end{insight}

\section{Riemannian metrics, lengths, and geodesics}
Topology alone says which deformations are continuous; it does not give a ruler. A Riemannian metric assigns a positive definite inner product $g_x$ to each tangent space, varying smoothly with $x$. In coordinates $u$, it is a positive definite matrix $G(u)$:
\begin{equation}
g_x(v,w)=a^TG(u)b,\qquad v=Dq(u)a,\quad w=Dq(u)b.
\end{equation}
For a manifold embedded in Euclidean space with the induced metric,
\begin{equation}\label{eq:induced-metric}
G(u)=Dq(u)^TDq(u).
\end{equation}
This follows because ambient squared length of a small displacement is $\norm{Dq(u)\,du}^2=du^TG(u)du$. Full-rank $Dq$ makes $G$ positive definite.

For a smooth coordinate curve $u(t)$, its length is
\begin{equation}
L[u]=\int_a^b\sqrt{\dot u(t)^TG(u(t))\dot u(t)}\,dt.
\end{equation}
The Riemannian distance between two points is the infimum of lengths of piecewise smooth paths joining them, within the same connected component. A geodesic solves the local stationary-length equations and, over sufficiently short segments, minimizes length. Long geodesic segments need not be globally minimizing. Existence of a shortest path between arbitrary points requires suitable completeness conditions.

On a circle of radius $R$, angular coordinate gives $G(\theta)=R^2$. If the smaller angular separation is $\alpha\in[0,\pi]$, intrinsic distance is $R\alpha$ whereas ambient chord distance is $2R\sin(\alpha/2)$. For small $\alpha$ these agree to leading order, motivating neighborhood approximations; for opposite points they differ by a factor $\pi/2$. On a sphere, great-circle arcs are geodesics; antipodal points admit multiple shortest arcs.

Coordinate change $u=u(v)$ transforms the metric by $G_v=(Du/Dv)^TG_u(Du/Dv)$. The numerical matrix changes while path lengths do not. This is why geometry is not the same as the numerical coordinates output by a manifold-learning algorithm.

\section{Curvature: bending within and outside a space}
For a regular plane curve parameterized by arc length $s$, curvature is $\kappa=\norm{dT/ds}$, where $T=dq/ds$ is its unit tangent. A circle of radius $R$ has curvature $1/R$. This measures how the curve bends in the plane. As a one-dimensional Riemannian manifold, a short arc has the same intrinsic length geometry as an interval; the ambient bending is extrinsic.

For surfaces, Gaussian curvature is the product of the two principal curvatures and is determined by the intrinsic metric. A sphere of radius $R$ has Gaussian curvature $1/R^2$. A cylinder has one nonzero principal curvature and one zero, so Gaussian curvature is zero: a small cylinder patch can be unrolled onto a plane without stretching even though it bends in three-dimensional space. The standard torus in \eqref{eq:torus-embedding} has regions of different Gaussian curvature signs. Its topology is independent of these particular curvature values.

This introductory distinction is enough for many neural interpretations. A trajectory bending in a PCA plot is extrinsic curvature of a chosen representation and projection. It does not estimate intrinsic Riemannian curvature without a metric, sufficient local sampling, and control of projection distortion. Likewise, two isometric patches can occupy very different shapes in ambient coordinates.

\section{What could justify a manifold model for data?}
The formal definition applies to a space, not directly to a finite noisy sample. A statistical model must connect observations to it, for example
\begin{equation}\label{eq:manifold-observation}
z_i\sim\rho\text{ on }M,\qquad y_i=g(z_i)+\eta_i,
\end{equation}
where $M$ is a $d$-manifold, $\rho$ a sampling distribution, $g$ an observation map, and $\eta_i$ observation noise. The manifold hypothesis is the assumption that such a low-dimensional description is useful on specified spatial, temporal, and noise scales. It is not an automatic consequence of fitting $d$ latent coordinates.

Several checks target different parts of the model. Stable local dimension estimates address the number of degrees of freedom. Consistent tangent estimates and local reconstruction errors address smoothness. Neighborhood preservation addresses the embedding. Sampling coverage addresses whether distant parts of the space have been visited. Topological methods address connectivity and holes. Out-of-sample prediction addresses whether the representation describes more than the training cloud. No one check establishes all these properties.

A ring-shaped projection is compatible with a circle, an annulus, a curved cylinder, an oscillating trajectory in a higher-dimensional state space, or clustered observations arranged by the algorithm. A mathematically precise claim must specify whether it concerns a latent manifold, the support of a distribution, an occupied subset, a fitted decoder image, or a visualization. In neural work, the operational meaning of ``manifold'' may intentionally be a useful low-dimensional subspace; the appropriate response is to state that meaning and test its claims, not to demand an unclaimed differential-geometric theorem.

\section{Exercises}
\begin{enumerate}
\item \textbf{Introductory.} Decide whether an open interval, closed interval, crossed pair of lines, and finite point cloud are manifolds without boundary. State the dimension and any exceptional points.
\item \textbf{Introductory.} Differentiate the circle parameterization and calculate its induced metric for radius $R$.
\item \textbf{Intermediate.} Verify the stereographic inverse and identify its missing point. Explain why a coordinate singularity need not be a singularity of the manifold.
\item \textbf{Intermediate.} Apply Theorem~\ref{thm:implicit-function} to the ellipsoid $x^2/a^2+y^2/b^2+z^2/c^2=1$ with positive $a,b,c$. Compute its tangent plane at a general point.
\item \textbf{Intermediate.} Compute the induced metric of \eqref{eq:torus-embedding}. Show that the two coordinate directions are orthogonal and that their lengths depend differently on position.
\item \textbf{Advanced.} Give a smooth observation map from a circle to a line that identifies distinct states. Explain how this violates the assumption required to recover the circle's topology from observations.
\item \textbf{Advanced.} Formulate a falsifiable manifold sampling hypothesis for neural activity, including dimension, noise, observation map, and sampling range. List separate tests for each assumption.
\end{enumerate}

\section{Summary and further reading}
A manifold is locally Euclidean with explicit topological regularity; a smooth manifold also has compatible differentiable charts. Tangent spaces describe local degrees of freedom, embeddings relate abstract states to measurements, and a Riemannian metric supplies lengths and angles. Topology, intrinsic geometry, and ambient shape are distinct. \citet{lee2012} gives a comprehensive mathematical development, while the derivations here provide the prerequisites needed for the operational neural comparisons in Chapter~\ref{ch:08}. Chapter~\ref{ch:04} turns global connectivity and holes into computable invariants.
